\chapter*{Guia de uso}
\addcontentsline{toc}{section}{Guia de uso}

\begin{theorybox}{Como esta organizado cada bloque}
Los capitulos de contenido siguen un patron estable: objetivos, prerrequisitos, teoria minima,
metodo explicito, ejemplo resuelto, error frecuente, ejercicios guiados, practica graduada,
autoevaluacion y, cuando procede, ampliacion.
\end{theorybox}

\begin{methodbox}{Ruta recomendada de trabajo}
\begin{enumerate}
  \item Empieza por la seccion o capitulo que mas necesites reforzar.
  \item Haz primero los guiados sin mirar la solucion completa.
  \item Usa las respuestas breves solo para comprobar si vas bien encaminado.
  \item Si fallas por eleccion de tecnica, pasa al bloque de repaso acumulativo.
  \item Reserva los simulacros para trabajo con tiempo y correccion posterior.
\end{enumerate}
\end{methodbox}

\begin{commonerror}{Confundir mirar con estudiar}
Leer una solucion correcta no equivale a haber ganado el metodo. Para consolidar un tipo de
ejercicio conviene cerrar la solucion, rehacer el problema y justificar cada paso con palabras
propias.
\end{commonerror}

\subsection*{Claves visuales}
\begin{itemize}
  \item Las cajas de \textbf{Teoria minima} resumen la idea imprescindible.
  \item Las cajas de \textbf{Metodo} explicitan el algoritmo de trabajo.
  \item Las cajas de \textbf{Error frecuente} senalan tropiezos reales que conviene anticipar.
  \item Las cajas de \textbf{Ampliacion} marcan contenido no nuclear o estilo EBAU.
  \item En la edicion del profesorado aparecen, ademas, las \textbf{soluciones completas}.
\end{itemize}

\subsection*{Cuando usar el bloque final}
El bloque de repaso acumulativo esta pensado para el tramo final:
\begin{itemize}
  \item diagnostico rapido al inicio del repaso;
  \item hoja mixta para recombinar tecnicas;
  \item seleccion de metodo cuando la dificultad no es operar, sino reconocer el enfoque;
  \item analisis de errores para entrenar vigilancia algebraica;
  \item simulacros con baremo y reto final de transferencia.
\end{itemize}
